\scanpage{35}
\mspage{22}

on a

\[
 u_{0,\alpha_0}\in C.
\]

Compte tenu du complément marginal à proposition 1 (qui se prouve de la $=$ façon, \unsure{\ldots}), on est ramené à prouver que \deleted{\unsure{des|les}} \inserted{l’hom. des} systèmes ordinaux envisagés est dans $C_{\unsure{\mathbf I|\mathbb I}_{\alpha_0+1}}$ \deleted{\unsure{a}} (où $\mathbf I_{\alpha_0+1}$ est l’ens des ordinaux $\leq\alpha_0$). Or \unsure{qu’il} en \unsure{soit ainsi} se \unsure{prouve} \inserted{\unsure{exact$^{\rm t}$|exact}} comme tantôt, à cela près \deleted{\unsure{qu’on}} \deleted{\unsure{que} \unsure{la} \unsure{vérific.}} \inserted{qu’on doit \inserted{\unsure{à présent}} faire} une vérification \unsure{du} cas p.~1 ($\widetilde u_{\unsure{v|\nu}_0}\in C$) pour $v_0=0$. Alors le diagramme \unsure{cas|de} p.~1 est

\diagram{\begin{tikzcd}[ampersand replacement=\&]
 \varnothing_{\deleted{\unsure{0|I}}}\arrow[r]\arrow[d,no head,"{\mathrm{iso}}"]\&X_0=X_{I'}\arrow[d,"{\mathrm{iso}}"',"{\mathrm{id}}"]\\
 \varnothing_{\deleted{\unsure{0|I}}}\arrow[r]\&X_0=X_{I'}\arrow[d,"{u_0=u_{I'}}"]\\
 \&Y_0=Y_{I'}
\end{tikzcd}}

et il faut donc vérifier $u_0=u_{I'}\in C$. Or ceci est contenu dans la prop.~2 elle-$=$, déjà prouvée, qu’on doit donc utiliser. Le reste se prouve comme avant. Ainsi, le cor.~2 de prop.~\unsure{2} est également prouvé

\uline{Commentaires} : Il y a un malaise, car les prop.~1 et prop.~2 \deleted{\unsure{sont}} (ainsi que leurs corollaires) sont \unsure{essentielt|essentiellem$^{\rm t}$} les $=$ \unsure{conclusions|conclusion}, avec des hypothèses \inserted{sur $I$} qui sont
% archive sheet number: 34
